\chapter{Technische Implementierung von Data-Masking}

\section{Verfahren der Anonymisierung}

Vor der detaillierten Betrachtung der einzelnen Verfahren bietet sich ein systematischer Vergleich der beiden zentralen Maskierungstechniken an:

\begin{table}[ht]
    \centering
    \footnotesize
    \caption{Vergleich der Maskierungstechniken}
    \label{tab:maskierungstechniken}
    \begin{tabular}{@{} >{\raggedright\arraybackslash}p{3.5cm} >{\raggedright\arraybackslash}p{5.5cm} >{\raggedright\arraybackslash}p{5.5cm} @{}}
        \toprule
        \textbf{Kriterium}    & \textbf{\ac{Regex}}                                 & \textbf{\ac{NER}}                            \\
        \midrule
        Erkennungsgenauigkeit & Hoch für strukturierte Daten (IP-Adressen, E-Mails) & Hoch für unstrukturierte, kontextuelle Daten \\
        Latenz                & Sehr gering (Millisekunden)                         & Höher (abhängig von Modellgröße)             \\
        Komplexität           & Niedrig (regelbasiert)                              & Hoch (KI-basiert, erfordert Training)        \\
        Wartbarkeit           & Manuelle Regelpflege erforderlich                   & Kontinuierliches Fine-Tuning nötig           \\
        Einsatzgebiet         & Formale, standardisierte Datenformate               & Freitext, variable Formulierungen            \\
        Ressourcenbedarf      & Minimal (CPU-basiert)                               & Hoch (GPU für Transformer-Modelle)           \\
        \bottomrule
    \end{tabular}
\end{table}

\subsection{Reguläre Ausdrücke}
\ac{Regex} gilt als verbreitetste Methode zur Identifikation von Mustern in Texten.
Es verwendet vordefinierte Zeichenketten, um spezifische Textmuster mit hoher Präzision in einer Zeichenkette zu finden.
Im Kontext von IT-Infrastrukturen können so z. B. IP-Adressen identifiziert werden und für die weitere Verarbeitung
mit anonymisierten Werten maskiert werden.

Der große Vorteil von \ac{Regex} ist dabei die schnelle und ressourcenschonende Verarbeitung der Daten, da moderne Regex-Engines hochoptimiert
für String-Verarbeitungsalgorithmen sind. Dabei können auch Datenbank-Verbindungsstrings, welche typischerweise Schlüsselwörter wie "server",
"database" oder "password" enthalten, zuverlässig erkannt und nachfolgende Informationen maskiert werden.
Die Implementierung einer solchen Lösung erfolgt typischerweise durch die Verwendung von \ac{Regex}-Bibliotheken oder externen Anbietern.

Die Limitation von \ac{Regex} liegt jedoch darin, dass es sich um eine rein regelbasierte Methode handelt.
So können z. B. Servernamen, welche einer bestimmten Namenskonvention folgen (z. B. "prod-db-server-01"), zuverlässig erkannt und maskiert werden.
Werden die Servernamen jedoch willkürlich gewählt (z. B. "der Hauptdatenbankserver"), kann \ac{Regex} diese nicht mehr erkennen,
da in diesem Fall semantische Zusammenhänge verwendet werden müssen und syntaktische Mustererkennung nicht ausreicht (Liu \& Vaidya, 2022).
Gleichzeitig besteht die Gefahr des Undermaskings: Neue oder unbekannte Namenskonventionen oder Datenformaten
werden nicht in die \ac{Regex}-Regeln aufgenommen und sommit nicht erkannt,
bzw. die Gefahr des Overmaskings, dass zu viele Daten maskiert werden und folglich die Antwortqualität der \ac{LLM}-\ac{API} sinkt.


\subsection{Named Entity Recognition}
Im Gegensatz zu \ac{Regex} adressiert \ac{NER} das Problem der Datenmaskierung mithilfe von semantischer Klassifikation.
Dabei wird ein kontextuelles Sprachverständnis genutzt, um die Bedeutung von Wörtern im umgebenden Text zu erfassen und somit auch
Entitäten zu erkennen, welche nicht explizit in Regeln definiert sind. \ac{NER} "denkt mit" und ein Servername wird nicht erkannt,
weil er aussieht wie ein Servername, sondern weil der Satzbau („...verbindet sich mit dem Host Alpha-Prime...“) darauf hindeutet.

Der Nachteil von \ac{NER} zeigt sich jedoch in der Komplexität und den Ressourcen, welche für das Training und den Betrieb benötigt werden.
So müssen die Modelle aufwendig trainiert und gewartet werden, was zeit- und kostenintensiv sein kann.
Aus diesem Grund lohnt sich der Einsatz von \ac{NER} nur, wenn eine hohe Genauigkeit bei der Erkennung von Entitäten erforderlich ist und die Ressourcen für das Training und den Betrieb vorhanden sind.

Als Alternative gibt es einen hybriden Kaskaden-Ansatz. Dabei wird die Anfrage zuerst mit minimaler Latenz
durch ein \ac{Regex}-basiertes Modell geschickt. Alles, was harte Muster enthält, wird erkannt und mit anonymisierten Namenskonventionen ersetzt.
Der verbleibende Text wird daraufhin an das \ac{NER}-basierte Sicherheitsnetz geschickt, welches verbleibende Informationen erkennt und maskiert.
Dieses System stellt dabei nicht nur eine hohe Effizienz bereit, sondern auch eine erhöhte Zuverlässigkeit, da sich die beiden Modelle ergänzen.
Um die anonymisierten Antworten der \ac{LLM}-\ac{API} wieder ersetzen zu können, ist es wichtig, einen gemeinsamen Token-Speicher zu nutzen, welcher
alle anonymisierten Begriffe speichert.


\section{Das Utility-Privacy-Dilemma}

Die zentrale Herausforderung bei einer erfolgreichen Datenmaskierung liegt in der Balance zwischen dem Schutz sensibler Daten (Privacy) und der
gleichzeitigen Erhaltung der Nützlichkeit der Daten (Utility) für die \ac{LLM}-\ac{API}. In der Wissenschaft
wird dies oft als Privacy-Utility-Trade-off (Dwork \& Roth, 2014) bezeichnet.
Maximale Antwortqualität für die \ac{LLM}-\ac{API} bedeutet maximale Datenmenge,
was wiederum maximale Risiken für die Privatsphäre bedeutet.
Umgekehrt bedeutet maximale Datenminimierung maximale Ungenauigkeit bei der Antwortqualität der \ac{LLM}-\ac{API}. So kann die \ac{LLM}-\ac{API} zum Beispiel
nicht mehr die Zusammenhänge der \ac{IT}-Infrastruktur verstehen und so fehlerhafte Antworten zurückgeben.
Als Lösung wird häufig ein „Label-Preserving Masking“ eingesetzt. Anstatt sensible Daten durch anonymisierte Werte (z. B. [Mask]) zu ersetzen,
werden diese durch Label (z.B. [IP-ADDRESS1]) ersetzt, welche im Kontext der Anfrage verstanden werden können. So bleibt die logische
Topologie der \ac{IT}-Infrastruktur erhalten und die Antwortqualität der \ac{LLM}-\ac{API} nicht beeinträchtigt.
